%% ============================================================
%%  PARTE 1 — Descripción del Liver Cancer Algorithm (LCA)
%%  Hasta 1.0 punto
%% ============================================================

\section{Parte I: El Liver Cancer Algorithm}
\label{sec:parte1}

\subsection{Motivación y contexto}

La búsqueda de nuevas metaheurísticas eficientes es una necesidad constante en el campo de
la optimización, especialmente ante el teorema \emph{No Free Lunch} \cite{wolpert1997nfl},
que establece que ningún algoritmo es universalmente superior a todos los demás para todo
tipo de problemas. En este contexto, el \textbf{Liver Cancer Algorithm (LCA)} surge como
una propuesta novedosa publicada en 2023 por Houssein \emph{et al.} \cite{houssein2023lca}
en la revista \emph{Computers in Biology and Medicine}.

La motivación del LCA reside en modelar matemáticamente el comportamiento agresivo y
adaptativo del \textbf{carcinoma hepatocelular} (cáncer de hígado primario). Esta enfermedad
presenta características biológicas especialmente ricas desde el punto de vista computacional:
el tumor crece de forma local y adaptativa, replica sus células en regiones cercanas y,
en estados avanzados, envía metástasis a órganos lejanos. Estas tres etapas se traducen
directamente en los tres operadores principales del algoritmo.

\subsection{Fundamentación biológica}

El LCA extrae su inspiración de los siguientes mecanismos oncológicos:

\begin{enumerate}[label=\textbf{\arabic*.}]
  \item \textbf{Crecimiento adaptativo del tumor} (\emph{explotación}): el tumor evalúa su
    entorno local y crece hacia las regiones con mayor disponibilidad de recursos (vascularización,
    nutrientes). En términos de optimización, esto equivale a una búsqueda local alrededor
    de la mejor solución conocida.

  \item \textbf{Replicación y proliferación celular} (\emph{transición}): las células tumorales
    se replican siguiendo un modelo de crecimiento exponencial. La tasa de crecimiento
    introduce un elemento estocástico controlado que permite explorar el vecindario inmediato
    con cierta amplitud.

  \item \textbf{Metástasis} (\emph{exploración global}): en la fase más avanzada, las células
    malignas viajan a través del torrente sanguíneo mediante saltos de tipo \emph{Lévy Flight},
    colonizando regiones lejanas e inesperadas del espacio de búsqueda.

  \item \textbf{Operadores genéticos de metástasis} (mutación y cruce): modelan la heterogeneidad
    genética del tumor, introduciendo variabilidad y combinando información de distintas
    soluciones para generar descendencia más apta.

  \item \textbf{ROBL — Random Opposition-Based Learning}: inicializa la población generando
    soluciones y sus opuestas aleatorias, asegurando una cobertura inicial diversa del
    espacio de búsqueda.
\end{enumerate}

\subsection{Modelo matemático detallado}

\subsubsection{Inicialización mediante ROBL}

Las posiciones iniciales de la población se generan mediante la ecuación:

\begin{equation}
  \label{eq:robl}
  \text{Position}^*_{j,i} = \frac{\pi}{6}(\ell_j)(w_j)(h_j)
    - \Bigl(lb + (ub - lb) - r_d \cdot \text{Position}^j_i\Bigr)
\end{equation}

donde $i = 1,\ldots,N$, $j = 1,\ldots,D$, $lb$ y $ub$ son los límites del espacio de
búsqueda, y $r_d \in [0,1]$ es un número aleatorio. Los parámetros $\ell, w, h$
(longitud, anchura, altura del tumor) son números aleatorios en $[0,1]$ que modelan las
dimensiones del hemi-elipsoide tumoral.\footnote{Un \emph{hemi-elipsoide} no
designa un tipo de tumor, sino un \textbf{modelo geométrico-matemático}: muchos
tumores sólidos (mama, hígado) se aproximan por una forma tridimensional de media
esfera achatada (medio elipsoide), de la que el LCA toma sus parámetros $\ell$, $w$
y $h$ para modelar el crecimiento.}

\subsubsection{Crecimiento adaptativo (explotación)}

El incremento de tamaño del tumor se calcula como:

\begin{equation}
  \label{eq:crec}
  \text{Position} = \frac{\pi}{6} \cdot f \cdot (l \cdot w)^{3/2}
\end{equation}

siendo $f=1$ una constante que caracteriza el tipo de tumor. Esta ecuación modela el
volumen del hemi-elipsoide y proporciona el desplazamiento local de cada agente de búsqueda.

\subsubsection{Replicación y vuelo de Lévy (exploración global)}

El crecimiento exponencial del tumor se rige por:

\begin{equation}
  \label{eq:repl}
  (PG)_i = \frac{dV}{dt} = r \cdot \text{Position},
  \quad t \in [1\ldots T],\; i \in [1\ldots N]
\end{equation}

La dispersión mediante vuelo de Lévy se implementa como:

\begin{equation}
  \label{eq:levy}
  L_v(D) = 0.01 \times \text{rand}(1,D) \times \frac{\sigma}{|\text{rand}(1,D)|^{1/\beta}}
\end{equation}

\begin{equation}
  \sigma = \left(\frac{\Gamma(1+\beta)\sin(\pi\beta/2)}
    {\Gamma\!\left(\frac{1+\beta}{2}\right)\cdot\beta\cdot 2^{(\beta-1)/2}}\right)^{1/\beta}
\end{equation}

La actualización de posición integra ambos mecanismos:

\begin{align}
  y &= \text{Position} + PG \label{eq:y}\\
  Z &= Y + S \cdot LF(D) \label{eq:Z}\\
  \text{Position}_i^{t+1} &=
    \begin{cases}
      y & \text{si } fit(y) < fit(\text{Position}_i^t) \\
      Z & \text{si } fit(Z) < fit(\text{Position}_i^t)
    \end{cases} \label{eq:update}
\end{align}

\subsubsection{Operadores genéticos: mutación y cruce}

La \textbf{mutación adaptativa} está controlada por el parámetro $\zeta = t/T$,
que evoluciona linealmente de 0 a 1 a lo largo de las iteraciones:

\begin{equation}
  y_{\text{Mut}} =
  \begin{cases}
    \text{Position} & \text{si } \rho_1 \geq \zeta \\
    y               & \text{en otro caso}
  \end{cases},
  \qquad
  z_{\text{Mut}} =
  \begin{cases}
    \text{Position} & \text{si } \rho_2 \geq \zeta \\
    z               & \text{en otro caso}
  \end{cases}
  \label{eq:mut}
\end{equation}

El \textbf{cruce} combina linealmente las dos versiones mutadas:

\begin{equation}
  w_{\text{Cross}} = \tau \cdot y_{\text{Mut}} + (1-\tau') \cdot z_{\text{Mut}},
  \quad \tau \neq \tau'
  \label{eq:cross}
\end{equation}

La \textbf{selección codiciosa} elige al mejor entre padre e hijo:

\begin{equation}
  \text{Position}_i^{t+1} =
  \begin{cases}
    y_{\text{Mut}}    & \text{si } fit(y_{\text{Mut}}) < fit(\text{Position}_i) \\
    z_{\text{Mut}}    & \text{si } fit(z_{\text{Mut}}) < fit(\text{Position}_i) \\
    w_{\text{Cross}}  & \text{si } fit(w_{\text{Cross}}) < fit(\text{Position}_i)
  \end{cases}
  \label{eq:sel}
\end{equation}

\subsection{Pseudocódigo del LCA}

\begin{algorithm}[H]
\caption{Liver Cancer Algorithm (LCA)}
\label{alg:lca}
\begin{algorithmic}[1]
\Require $N$ (tamaño de población), $T$ (iteraciones máx.), $lb$, $ub$ (límites)
\Ensure $X_{\text{best}}$ (mejor solución encontrada)
\State Inicializar posiciones $\{\text{Position}_i\}_{i=1}^N$ mediante ROBL (Eq.~\ref{eq:robl})
\State Evaluar fitness de cada agente $i$
\State Identificar $X_{\text{best}}$
\For{$t = 1$ \textbf{to} $T$}
  \State $\zeta \leftarrow t / T$ \Comment{tasa de mutación adaptativa}
  \For{cada agente $i \in [1, N]$}
    \State Calcular volumen de crecimiento $V$ y $PG = r\cdot V$ (Eqs.~\ref{eq:crec}--\ref{eq:repl})
    \State $y \leftarrow \text{Position}_i + PG$ \Comment{crecimiento adaptativo (\textit{explotación})}
    \State $Z \leftarrow y + S\cdot LF(D)$ \Comment{vuelo de Lévy (\textit{exploración})}
    \State Seleccionar entre $y$ y $Z$ según fitness (Eq.~\ref{eq:update})
    \State Aplicar mutación con tasa $\zeta$ (Eq.~\ref{eq:mut})
    \State Aplicar cruce (Eq.~\ref{eq:cross})
    \State Selección codiciosa (Eq.~\ref{eq:sel})
    \State Evaluar nuevo fitness
    \If{$fit(\text{Position}_i^{t+1}) < fit(X_{\text{best}})$}
      \State $X_{\text{best}} \leftarrow \text{Position}_i^{t+1}$
    \EndIf
  \EndFor
\EndFor
\State \Return $X_{\text{best}}$
\end{algorithmic}
\end{algorithm}

\paragraph{Nota sobre la fidelidad del pseudocódigo.}
El pseudocódigo del Algoritmo~\ref{alg:lca} es una \textbf{reconstrucción coherente}
elaborada a partir de las ecuaciones del artículo original (Eqs.~\ref{eq:robl}--\ref{eq:sel}),
que constituyen la especificación matemática real del método. No reproduce literalmente el
bloque ``Algorithm~1'' de Houssein \emph{et al.} \cite{houssein2023lca}, ya que este presenta
varias inconsistencias formales: incluye una condición tautológica
(\texttt{if (Position $\le 0.8$) or (Position $>0.8$)}, cierta para cualquier valor), compara
el vector de posición $D$-dimensional con umbrales escalares ($0.8$, $1.96$) sin precisar sobre
qué componente, y combina ramas \texttt{if}/\texttt{else if} de forma sintácticamente ambigua.
Por ello se ha optado por la versión secuencial bien formada ---crecimiento, vuelo de Lévy,
selección codiciosa y operadores genéticos--- que el propio texto del artículo (secciones~3.1.2
y~3.1.3) describe y respaldan sus ecuaciones. Las diferencias entre este pseudocódigo y la
implementación efectivamente evaluada se detallan en la Parte~II
(\S\,\ref{sec:parte2}, ``Desviaciones respecto a la formulación original'').

\subsection{Complejidad temporal}

La complejidad total del LCA es $\mathcal{O}(N \cdot (T + T \cdot D + 1))$, donde $N$
es el tamaño de la población, $T$ el número de iteraciones y $D$ la dimensión. Las tres
fases que dominan el coste son: inicialización $\mathcal{O}(N)$, evaluación de fitness
$\mathcal{O}(T \cdot N)$ y actualización de posiciones $\mathcal{O}(T \cdot N \cdot D)$.
En la práctica, con $N=30$ y los límites de evaluaciones del CEC2017, este coste es perfectamente
asumible en tiempos de ejecución razonables.

\subsection{Análisis del equilibrio exploración–explotación}

Uno de los criterios de evaluación más importantes de la práctica es el análisis del
equilibrio entre \textbf{exploración} e \textbf{explotación}. La tabla siguiente sintetiza
cómo cada componente del LCA contribuye a uno u otro aspecto:

\begin{table}[h!]
\centering
\caption{Componentes del LCA y su rol en el equilibrio exploración–explotación}
\label{tab:equilibrio}
\renewcommand{\arraystretch}{1.35}
\footnotesize
\begin{tabular}{@{}lll@{}}
\toprule
\textbf{Componente} & \textbf{Mecanismo matemático} & \textbf{Rol principal} \\
\midrule
ROBL (inicialización) & Eq.~\ref{eq:robl} — soluciones opuestas & Diversidad inicial \\
Crecimiento adaptativo & Eq.~\ref{eq:crec} — volumen hemi-elipsoide & \textbf{Explotación} local \\
Replicación exponencial & Eq.~\ref{eq:repl} — $dV/dt = r \cdot \text{Pos}$ & Transición suave \\
Vuelo de Lévy & Eqs.~\ref{eq:levy}--\ref{eq:update} — saltos largos & \textbf{Exploración} global \\
Mutación adaptativa & Eq.~\ref{eq:mut} — $\zeta = t/T$ crece con $t$ & Exploración $\to$ Explotación \\
Cruce lineal & Eq.~\ref{eq:cross} — combinación de vectores & Combina ambas fases \\
Selección codiciosa & Eq.~\ref{eq:sel} — mejor padre/hijo & Elitismo (explotación) \\
\bottomrule
\end{tabular}
\end{table}

\paragraph{Discusión del equilibrio.}
El LCA logra un balance dinámico y progresivo entre exploración y explotación a través
del parámetro $\zeta = t/T$. Al inicio de la búsqueda ($t \approx 0$), $\zeta \approx 0$,
por lo que la mutación actúa con poca intensidad y predomina la exploración mediante el
vuelo de Lévy. A medida que avanza la ejecución ($t \to T$), $\zeta \to 1$ aumenta la
tasa de mutación y cruce, intensificando la explotación en torno a las mejores regiones
encontradas. En la formulación original, cada iteración evalúa \textbf{ambos} movimientos
---el crecimiento local $y = \text{Position}+PG$ (explotación) y el salto de Lévy
$Z = y + S\cdot LF$ (exploración)--- y conserva el de mejor fitness mediante selección
codiciosa (Eq.~\ref{eq:update}), combinando así exploración y explotación en cada paso.
Este diseño evita tanto la convergencia prematura —frecuente en algoritmos puramente
explotadores— como la exploración excesiva sin refinar las soluciones prometedoras.
